\SourceLeaf{72}
\Needspace{14\baselineskip}\begin{center}\textbf{[BEWEIS EINES VON EULER AUFGESTELLTEN SATZES ÜBER EXAKTE DIFFERENTIALAUSDRÜCKE.]}\end{center}



\noindent\hspace*{\parindent}Nimmt man $dW = Vdx$ an, wobei $W$ eine Funktion von $x, y, p, q, r, s,$ bezeichnet, so gilt

\[
\begin{aligned}V ={}&\left(\frac{\partial W}{\partial x}\right) + p\left(\frac{\partial W}{\partial y}\right) + q\left(\frac{\partial W}{\partial p}\right) \\ &{}+ r\left(\frac{\partial W}{\partial q}\right) + s\left(\frac{\partial W}{\partial r}\right) + t\left(\frac{\partial W}{\partial s}\right).
\end{aligned}
\]

\par\noindent Daraus folgt

\[
N = \left(\frac{\partial V}{\partial y}\right) = \frac{d\left(\frac{\partial W}{\partial y}\right)}{dx},
\]
\[
P = \left(\frac{\partial V}{\partial p}\right) = \left(\frac{\partial W}{\partial y}\right) + \frac{d\left(\frac{\partial W}{\partial p}\right)}{dx},
\]
\[
Q = \left(\frac{\partial V}{\partial q}\right) = \left(\frac{\partial W}{\partial p}\right) + \frac{d\left(\frac{\partial W}{\partial q}\right)}{dx},
\]
\[
R = \left(\frac{\partial V}{\partial r}\right) = \left(\frac{\partial W}{\partial q}\right) + \frac{d\left(\frac{\partial W}{\partial r}\right)}{dx},
\]
\[
S = \left(\frac{\partial V}{\partial s}\right) = \left(\frac{\partial W}{\partial r}\right) + \frac{d\left(\frac{\partial W}{\partial s}\right)}{dx},
\]
\[
T = \left(\frac{\partial V}{\partial t}\right) = \left(\frac{\partial W}{\partial s}\right).
\]

